\documentclass[12pt]{article}
\usepackage{amsmath,amssymb}

\begin{document}
\section*{{2022 AIME II Problems/Problem 2}}
Source: AoPS Wiki

\subsection*{Solution}
Let $A$ be Azar, $C$ be Carl, $J$ be Jon, and $S$ be Sergey. The $4$ circles represent the $4$ players, and the arrow is from the winner to the loser with the winning probability as the label. This problem can be solved by using $2$ cases. $\textbf{Case 1:}$ $C$ 's opponent for the semifinal is $A$ The probability $C$ 's opponent is $A$ is $\frac13$ . Therefore the probability $C$ wins the semifinal in this case is $\frac13 \cdot \frac13$ . The other semifinal game is played between $J$ and $S$ , it doesn't matter who wins because $C$ has the same probability of winning either one. The probability of $C$ winning in the final is $\frac34$ , so the probability of $C$ winning the tournament in case 1 is $\frac13 \cdot \frac13 \cdot \frac34$ $\textbf{Case 2:}$ $C$ 's opponent for the semifinal is $J$ or $S$ It doesn't matter if $C$ 's opponent is $J$ or $S$ because $C$ has the same probability of winning either one. The probability $C$ 's opponent is $J$ or $S$ is $\frac23$ . Therefore the probability $C$ wins the semifinal in this case is $\frac23 \cdot \frac34$ . The other semifinal game is played between $A$ and $J$ or $S$ . In this case it matters who wins in the other semifinal game because the probability of $C$ winning $A$ and $J$ or $S$ is different. $\textbf{Case 2.1:}$ $C$ 's opponent for the final is $A$ For this to happen, $A$ must have won $J$ or $S$ in the semifinal, the probability is $\frac34$ . Therefore, the probability that $C$ won $A$ in the final is $\frac34 \cdot \frac13$ . $\textbf{Case 2.2:}$ $C$ 's opponent for the final is $J$ or $S$ For this to happen, $J$ or $S$ must have won $A$ in the semifinal, the probability is $\frac14$ . Therefore, the probability that $C$ won $J$ or $S$ in the final is $\frac14 \cdot \frac34$ . In Case 2 the probability of $C$ winning the tournament is $\frac23 \cdot \frac34 \cdot (\frac34 \cdot \frac13 + \frac14 \cdot \frac34)$ Adding case 1 and case 2 together we get $\frac13 \cdot \frac13 \cdot \frac34 + \frac23 \cdot \frac34 \cdot (\frac34 \cdot \frac13 + \frac14 \cdot \frac34) = \frac{29}{96},$ so the answer is $29 + 96 = \boxed{\textbf{125}}$ . ~ isabelchen

\end{document}
